%قبل از شروع مسئله ترجیحاً چیزی ننویسید. اما اگر خواستید بنویسید، در قالب لیستی (فقرات و شمارش) بنویسید.
%برای اضافه کردن نمره، بعدی از شروع‌مسئله، [x] را اضافه کنید و به جای x  نمره مورد نظر را قرار دهید. برای زیرمسئله‌ها هم به همین صورت است.
%\شروع‌مسئله[10]
\شروع‌مسئله
% صورت اصلی مسئله را در این‌جا بنویسید.
با استفاده از داده ساختار پشته به سوالات زیر پاسخ دهید.

%اگر مسئله چند قسمت داشت به این صورت عمل کنید:
\شروع{زیرمسائل}
\زیرمسئله
با استفاده از پشته، داده ساختار صف را پیاده‌سازی کنید به طوری که هزینه عملیات اضافه کردن به صف در آن از $O(1)$ باشد. هزینه سرشکن $n$ عمل اضافه و حذف از صف را با استفاد از روش تابع پتانسیل تحلیل کنید. فرض کنید هزینه ها معادل تعداد $push$ها و $pop$ها است.
\زیرمسئله
به کمک پشته و حافظه اضافی از مرتبه $O(1)$ یک داده ساختار جدید مانند پشته ارائه دهید که بتواند علاوه بر عملیات‌های $push$ و $pop$، عملیات‌های $find\_min$ و $find\_max$ را از مرتبه زمانی $O(1)$ انجام دهد.
\زیرمسئله
آیا می‌توان داده ‌ساختار قسمت ب را به گونه‌ای تغییر دهیم که عملیات $delete\_min$ را نیز در  $O(1)$  انجام دهد؟
\پایان{زیرمسائل}
\solution{
	می‌دانیم که پشته‌ها FILO و صف‌ها FIFO هستند. حال فرض کنید که دو پشته داریم که پشته‌ی اول پر و پشته‌ی دوم خالی می‌باشد. حال به ترتیب اجزای پشته‌ی اول را خوانده و در پشته‌ی دوم می‌ریزیم، چون پشته‌ها FILO هستند، اجزایی که زودتر وارد پشته‌ی اول شده‌اند دیرتر وارد پشته‌ی دوم می‌شوند. چون پشته‌ی دوم نیز FILO است، اجزایی که دیرتر وارد شده‌اند زودتر خارج می‌شوند، پس هر عضوی که زودتر وارد پشته‌ی اول شده باشد، بعد از خالی کردن پشته‌ی اول به پشته‌ی دوم زودتر خارج می‌شود که مانند یک FIFO می باشد. حال با کمک این ایده به کمک دو پشته یک صف را طراحی می‌کنیم:
	
	هنگام push به پشته‌ی اول push می‌کنیم. هنگام pop نیز از پشته‌ی دوم pop میکنیم، اما اگر پشته‌ی دوم خالی بود، ابتدا تمام اعضای پشته‌ی اول را وارد پشته‌ی دوم می‌کنیم. سپس از پشته‌ی دوم pop میکنیم. عملیات push چون عملیات push به پشته است از 
	$O(1)$
	است. عملیات pop شامل حداکثر $n$ عمل pop برای خالی کردن پشته‌ی اول $n$ عمل push برای وارد شدن به پشته‌ی دوم و یک عملیات pop از پشته‌ی دوم می‌باشد. پس از
	$O(n)$
	می‌باشد. (گر چه پس از یادگیری تحلیل سرشکن می‌توانید بررسی کنید که مرتبه‌ی زمانی pop نیز به طور سرشکن از
	$O(1)$
می‌باشد چون هر عضو حداکثر دو بار pop می‌شود.)

برای محاسبه هزینه سرشکن:

$\phi=2len(S_{A})$


$\hat c_{i} = c_{i} + \phi_{i} - \phi_{i-1}$


چون در این سوال هزینه ی push و pop یکسان است پس داریم :

$c_{i pop} = c_{i push}$


برای Enqueue فقط یک push در استک A داریم که $c_{i}$ در آن یک است که به علاوه ی دوبرابر تفاوت فاصله ی استک A قبل و بعد از push کردن می‌شود.


$Enqueue \rightarrow \hat c_{i} = 1 + 2(len(S_{A})-len(S_{A-1}))  = 1+2= 3 = O(1)$


برای Dequeue ابتدا حالتی را در نظر میگیریم که استک B خالی نباشد. در این صورت تنها کافی است یکی از B برداریم.


$Dequeue \rightarrow if len(S_{B})>0 \rightarrow \hat c_{i} = 1 + 2(len(S_{A}) - len(S_{A-1})) = 1 + 0 = 1 = O(1)$


در حالتی که استک B خالی است باید تمام استک A را پاپ کنیم و بعد به اندازه ی طول A در استک B بریزیم و در آخر یکی pop می‌کنیم.


$c_{i} = 2len(S_{A}) + 1$


$\phi_{i} = 0$


$\phi_{i-1} = 2len(S_{A})$


بنابراین
$\hat c_{i}  = 1 = O(1)$



بنابراین هزینه سرشکن برای این پیاده سازی برای Enqueue برابر 3 است و برای Dequeue برابر 1 است.

قسمت ب و پ:
داده ساختار اصلی یک پشته برای عملیات‌های push و pop است. برای یافتن min نیز از یک پشته‌ی کمکی هم ارتفاع با پشته‌ی اصلی استفاده می‌کنیم که در ارتفاع $h$م آن، کمینه‌ی $h$ عنصر اول پشته‌ی اصلی را نگه می‌داریم. برای مدیریت این استک نیز به صورت زیر عمل می‌کنیم:

هنگام pop از استک اصلی، از این استک نیز pop می‌کنیم.
هنگام push کردن به استک اصلی، عضو جدید را با اخرین عضو پشته  کمکی مقایسه می‌کنیم. اگر کوچکتر بود، عضو جدید وگرنه آخرین عضو پشته‌ی کمکی را به پشته‌ی کمکی push می‌کنیم. از آنجایی که عضو انتهایی کمینه‌ی سایر اعضا تا به اینجا بوده‌است، با روشی که ارائه کردیم، عضو جدید پشته‌ی کمکی، کمینه‌ی اعضای پشته‌ی اصلی می‌باشد.
برای عملیات پیدا کردن ماکسیمم نیز یک استک اضافه به همین شکل میگیریم.
برای
\متن‌لاتین{deleteMin}،
نمی‌توان این کار را انجام داد. برای اثبات فرض کنید که بشود deleteMin را در مرتبه‌ی $O(1)$ انجام داد.
در این صورت تمام اعضای یک ارایه را در $O(n)$ درون این داده ساختار می‌ریزیم سپس یکی یکی کمینه را یافته در یک آرایه‌ی جدید ریخته و کمینه را از داده‌ساختار فعلی حذف کند تا داده ساختار خالی شود و هر کدام از این عملیات‌ها در $O(1)$ و در کل در $O(n)$ انجام می‌شود. در این صورت
توانسته‌ایم که یک آرایه‌ی دلخواه را در زمان $O(n)$ مرتب کنیم که طبق قضایا می‌دانیم که مرتب‌سازی آرایه حداقل از مرتبه‌ی زمانی
$\Omega(n\log(n))$
می‌باشد و این تناقض می‌باشد.

}



%%%%%%%%%%%%%%%%%%%%%%%%%%%%%%%%%%%%%%%%%%%%%%%%%%%%%
% برخی دستورات مفید برای تمرین‌ها. %
%%%%%%%%%%%%%%%%%%%%%%%%%%%%%%%%%%%%%%%%%%%%%%%%%%%%%

% فرمول‌های نمونه:
%$f(x) = \Theta(x)$
%$T(x, y) \leq \log(x) + 2^y + \binom{2n}{n}$
% {\tt code\_function}


% لیست موارد بدون شماره‌گذاری:
%\شروع{فقرات}
%    \فقره مورد اول 
%        \شروع{فقرات}
%            \فقره مورد اول داخلی 
%                \شروع{فقرات}
%                    \فقره اولین مورد خیلی داخلی 
%                    \فقره دومین مورد خیلی داخلی 
%                \پایان{فقرات}
%            \فقره مورد دوم داخلی 
%        \پایان{فقرات}
%    \فقره مورد دوم 
%\پایان{فقرات}

% لیست موارد شماره‌گذاری شده:
%\شروع{شمارش}
%    \فقره یک مورد 
%    \فقره یک مورد دیگر 
%    \فقره باز هم یک مورد دیگر
%\پایان{شمارش}

% برای نوشتن یک فرمول چند خطی که یک شماره بخوره. برای حالت‌های دیگر به سایت زیر مراجعه کنید:
% https://www.overleaf.com/learn/latex/Aligning_equations_with_amsmath
%\begin{align} 
%    \begin{split}
%        (x+y)^3 &= (x+y)^2(x+y) \\
%                &= (x^2+2xy+y^2)(x+y) \\
%                &= (x^3+2x^2y+xy^2) + (x^2y+2xy^2+y^3) \\
%                &= x^3+3x^2y+3xy^2+y^3
%    \end{split}                                 
%\end{align}

% نوشتن ماتریس و آرایه دوبعدی
%\begin{align}
%    A = 
%    \begin{bmatrix}
%        A_{11} & A_{21} \\
%        A_{21} & A_{22}
%    \end{bmatrix}
%\end{align}

% اضافه کردن تصویر با شرح. برای اضافه کردن تصویر بدون شرح، فقط دستور درج‌تصویر را استفاده کنید.
%\شروع{شکل}[h!]
%\تنظیم‌ازوسط
%  \درج‌تصویر[پهنا=0.5\پهنای‌متن]{img.jpg}
%\شرح{شرح تصویر}
%\برچسب{تصویر:۱}
%\پایان{شکل}
%\پایان{زیرمسائل}

%کد زیر را ببینید.
%%\lstinputlisting{source_filename.py}
%\begin{latin}
%\begin{lstlisting}
%# Import the modules
%import sys
%import random
%
%ans = True
%
%while ans:
%question = raw_input("Ask the magic 8 ball a question: (press enter to quit) ")
%
%answers = random.randint(1,8)
%
%if question == "":
%sys.exit()
%
%elif answers == 1:
%print "It is certain"
%
%elif answers == 2:
%print "Outlook good"
%
%elif answers == 3:
%print "You may rely on it"
%
%elif answers == 4:
%print "Ask again later"
%
%elif answers == 5:
%print "Concentrate and ask again"
%
%elif answers == 6:
%print "Reply hazy, try again"
%
%elif answers == 7:
%print "My reply is no"
%
%elif answers == 8:
%print "My sources say no"
%\end{lstlisting}
%\end{latin}
